\section{The full compensated-minor correction}
\label{sec:compensatedjet}
Fix \(\delta\in(0,1)\) and \(K_0>0\).
Let \(D=J_N/N+B\) satisfy Theorem~\ref{thm:secondpermanent}
with entry parameter \(A\leq(\log N)^{K_0}\).
Use its block matrices \(\Omega,\Gamma\), and set
\[
 P=\begin{pmatrix}0&D\\D^\top&0\end{pmatrix},\quad
 c=\frac{(\one,\one)}{\sqrt{2N}},\quad
 s=\frac{(\one,-\one)}{\sqrt{2N}},\quad
 K=(I+\Omega)^{-1}-\tfrac12cc^\top-ss^\top .
\]
Thus \(K\) is the inverse of \(I+P\) on \(s^\perp\),
extended by zero on \(s\).
Define \(h_i=-\Gamma_{ii}\) and
\[
 v=(\tfrac12I-K)h .
\]
The two halves of \(h\) have equal sums.
All these matrices are real, and
\[
 \max|\Gamma_{ij}|\leq C_\delta A/N,\quad
 \|K\|_{\infty\to\infty}\leq C_\delta(1+A),\quad
 \|v\|_\infty\leq C_\delta(1+A)^2/N .
\]

\begin{theorem}[A uniform signed first-order minor expansion]
\label{thm:minorjet}
For arbitrary row and column sets \(I,J\) with
\(|I|=|J|=r\leq(\log(2N))^3\), put
\(u=(\one_I,\one_J)\), retaining their separate coordinates.
There is a fixed finite exponent \(M=M(\delta,K_0)\) such that,
uniformly in these sets,
\[
 \left(\frac{\operatorname{per}D[I^c,J^c]}
                   {\operatorname{per}D}\right)^2
 =1+\frac{r(r-1)}N+u^\top\Gamma u+2v^\top u
                         +O_{\delta,K_0}\left(
                              \frac{\log(2N)^M}{N^2}\right).
\]
The two-support quadratic term and the compensation field
are part of the expansion. They are not assumed to vanish.
\end{theorem}
\begin{proof}
Write \(n=N-r\), \(Q=D[I^c,J^c]\), and
\(s_{IJ}=\sum_{i\in I,j\in J}D_{ij}\).
Let \(\ell=Pu=(D\one_J,D^\top\one_I)\).
The remaining row defects are the corresponding coordinates
of \(\ell\), each at most \(rA/N\), and their sum is
\(2(r-s_{IJ})\).
Let \(\widetilde D=\diag(1+z_A)Q\diag(1+z_B)\)
be its exact stochastic repair, with \(z=0\) on removed
coordinates and gauge \(s^\top z=0\).
The established repair gives
\(\|z\|_\infty\leq C_\delta rA\log(2N)/N\).

We give the local error accounting needed for a two-power
remainder. On retained coordinates the exact row equation
is
\[
 (I+P)z=\ell+\ell\mathbin{\odot}z
                         -z\mathbin{\odot}(Pz).
\]
On removed coordinates the discrepancy is \((Pz)-\ell\).
Its magnitude is at most a polynomial in \(r,A,\log(2N)\)
divided by \(N\).
Every off-diagonal entry of \(K\) is \(O_\delta((1+A)/N)\).
Multiplying that discrepancy, supported on \(2r\) coordinates,
by \(K\), and the retained quadratic discrepancy by the
row-sum bound for \(K\), proves
\[
 z=K\ell+
 O_\delta\left(
  \frac{\operatorname{poly}(r,1+A,\log(2N))}{N^2}\right)
 \quad\hbox{on retained coordinates}.
\]
Here and below \(\operatorname{poly}\) denotes a fixed finite
polynomial with nonnegative coefficients depending only on
the displayed fixed parameters. No degree grows with \(N,r,A\).
On removed coordinates \(K\ell\) can instead be bounded by
\(C_\delta rA\log(2N)/N\), which is multiplied by an
additional \(O_\delta(A/N)\) factor whenever used below.

Sum the exact nonlinear scaling equations and expand
\(\log(1+z)\) through degree two. With the actual retained
row degrees \(d_i\), they give
\[
 \sum d_i z_i=r-s_{IJ}-\tfrac12z^\top Q_{\rm bip}z,
\]
and hence the full scaling cost is
\[
 {\cal S}:=\sum\log(1+z_A)+\sum\log(1+z_B)
 =r-s_{IJ}+\tfrac12\ell^\top K\ell
       +O_\delta\left(
       \frac{\operatorname{poly}(r,1+A,\log(2N))}{N^2}\right).
\]
Indeed the cubic Taylor cost is at most \(2N\|z\|_\infty^3\);
the omitted-coordinate costs are at most
\(2r(rA/N)\|K\ell\|_\infty\).
Both have the stated scale. The retained-coordinate approximation and
\(\|\ell\|_1=2r\) control the remaining substitutions.
This step retains the quadratic compensation, not just
the maximum scaling multiplier.

Next we compute the determinant cost locally.
Put \(M=I+\Omega\), and let \(S\) be the \(2r\) removed
coordinates. Jacobi's identity gives
\[
 \frac{\det M[S^c]}{\det M}
     =\det(M^{-1}[S])
     =\det(I-\Gamma[S]).
\]
Since \(\|\Gamma[S]\|\leq C_\delta rA/N=o(1)\),
its contribution to the logarithm of the inverse-square-root
determinant is
\[
 \tfrac12\operatorname{tr}\Gamma[S]
       +O_\delta(r^2A^2/N^2)=-\tfrac12 h^\top u+
                                      O_\delta(r^2A^2/N^2).
\]
The remaining changes are scaling and the change
\(J_n/N\) to \(J_n/n\).
The linear scaling contribution to the log determinant is
\(h^\top K\ell\).
For detail the off-diagonal block of \(M^{-1}\) is
\(-B(I-B^\top B)^{-1}\); its row and column sums vanish.
Contracting it with \(\diag(z_A)D+D\diag(z_B)\)
therefore gives \(h^\top z\), with the constant \(J_N/N\)
parts cancelling exactly.

All other determinant terms have a two-power remainder.
The Schur-complement change from \(M^{-1}[S^c]\) to
\((M[S^c])^{-1}\) has trace norm at most \(C_\delta rA/N\).
It is multiplied by a perturbation of operator norm
\(O_\delta(rA\log(2N)/N)\).
The squared Frobenius norm of the scaling perturbation is
\(O_\delta(r^2(1+A)^3\log(2N)^2/N^2)\).
Removed-coordinate contractions contain \(rN\) terms,
each with one scaling factor and two \(O_\delta(A/N)\)
matrix entries.
Finally the constant-block change contracts to zero on the
full supports; its remaining restricted contraction contains
only the removed row/column overlap and is of the same
smaller scale. The exact quadratic log-determinant remainder
is bounded by these squared Frobenius norms, since the
matrices retain a fixed positive spectral lower bound.
Consequently, writing \(\widetilde B=\widetilde D-J_n/n\),
\[
 \log\frac{\det(I-\widetilde B^\top\widetilde B)^{-1/2}}
               {\det(I-B^\top B)^{-1/2}}
 =-\tfrac12h^\top u+h^\top K\ell
 +O_\delta\left(
       \frac{\operatorname{poly}(r,1+A,\log(2N))}{N^2}\right).
\]

The first permanent correction must also be compared,
rather than independently discarded at order \(N^{-1}\).
The same Schur formula and scaling estimates imply, on
retained coordinates,
\[
 |\widetilde\Gamma_{ij}-\Gamma_{ij}|
 \leq\frac{\operatorname{poly}(r,1+A,\log(2N))}{N^2}.
\]
To check this entry bound, both full inverse matrices have
off-diagonal entries \(O_\delta(A/N)\) and absolute row
sums \(O_\delta(1+A)\). The entrywise inverse perturbation
identity therefore preserves the displayed scale.
Deleted coordinates contribute at most \(2r\) diagonal
terms or \(4rN\) off-diagonal terms.

Apply this estimate to every displayed term of \(\Psi\):
a diagonal square has difference bounded by its size
\(O(A/N)\) times the entry difference; the cubic sums have
two such small factors.
After summation their differences are
\(\operatorname{poly}(r,1+A,\log(2N))/N^2\).
For the Gamma averaging term, set
\[
 S_1=\operatorname{tr}\frac{B^\top B}{I-B^\top B},\qquad
 S_2=\operatorname{tr}
                      \left(\frac{B^\top B}{I-B^\top B}\right)^2.
\]
Then \(\Delta''(1)/\Delta(1)=S_1^2+S_1+2S_2\).
Here \(S_1\) is minus half the trace of \(\Gamma\), and
\(S_2\) is half the sum of squared entries of its two
diagonal blocks.
The preceding local bounds change each by
\(\operatorname{poly}(r,1+A,\log(2N))/N\).
The additional factor \(1/(N-1)\), including its change to
\(1/(n-1)\), proves
\[
 {\cal Q}_n(\widetilde D)-{\cal Q}_N(D)
       =O_\delta\left(
       \frac{\operatorname{poly}(r,1+A,\log(2N))}{N^2}\right).
\]
This is why the explicit first correction in the full
permanent theorem was needed.

Now use Theorem~\ref{thm:secondpermanent} in numerator
and denominator, together with the exact scaling cost.
Stirling gives
\[
 \log\frac{n!/n^n}{N!/N^N}
       =r-\frac r{2N}+O((r+r^2)/N^2).
\]
The scaling, determinant, and permanent-correction estimates yield
\[
 \log\frac{\operatorname{per}Q}{\operatorname{per}D}
 =s_{IJ}-\tfrac12\ell^\top K\ell-\frac r{2N}
           +(\tfrac12h-Kh)^\top u
           +O_{\delta,K_0}(\log(2N)^M/N^2).
\]
Finally \(u\perp s\), \(\ell=Pu\), and on \(s^\perp\)
\(PKP=P-I+K\). Since \(u^\top Pu=2s_{IJ}\),
\(u^\top u=2r\), and
\[
 u^\top Ku=2r-u^\top\Gamma u-r^2/N,
\]
the scalar and overlap terms simplify to
\[
 \frac{r(r-1)}{2N}+\tfrac12u^\top\Gamma u+v^\top u.
\]
This quantity is a fixed polynomial in \(r,A,\log(2N)\)
divided by \(N\). Exponentiating twice changes it only by
a two-power remainder and proves the displayed expansion.
Every polynomial above becomes a fixed power of \(\log(2N)\)
on \(r\leq\log(2N)^3\), \(A\leq(\log N)^{K_0}\).
\end{proof}


\begin{remark}[Independent check against every cofactor recursion]
Let \(W(I,J)=r(r-1)/(2N)+u^\top\Gamma u/2+v^\top u\)
be the first-order logarithmic correction, and put \(y=Kh\).
For \(i\notin I\), \(j\notin J\), its exact increment is
\[
 \begin{split}
 W(I\cup\{i\},J\cup\{j\})-W(I,J)
 ={}&r/N+(\Gamma u)_{A_i}+(\Gamma u)_{B_j}\\
    &+\Gamma_{A_i,B_j}-y_{A_i}-y_{B_j}.
 \end{split}
\]
Extend this polynomial expression temporarily to every \(j\),
only for the algebraic summation.
The exact matrix identities
\[
 (I+P)\Gamma=\Omega,\qquad (I+P)y=h,\qquad
 \sum_jD_{ij}\Gamma_{A_i,B_j}=h_{A_i}
\]
show that its \(D_{ij}\)-weighted sum over all \(j\) is
\[
 r/N+\bigl((I+P)\Gamma u\bigr)_{A_i}
             =r/N+(B\one_J)_i=(D\one_J)_i .
\]
Deleting the \(r\) forbidden \(j\)'s changes the sum by
only \(\operatorname{poly}(r,1+A,\log(2N))/N^2\).
Multiplying \(1+W(I,J)\) by the actual remaining row
mass \(1-(D\one_J)_i\) proves
\[
 \sum_{j\notin J}D_{ij}
       \bigl(1+W(I\cup\{i\},J\cup\{j\})\bigr)
 =1+W(I,J)+O_{\delta,K_0}(\log(2N)^M/N^2).
\]
This is precisely the first-order form of the exact
permanent cofactor recursion, for every support pair,
not just its row average. The column check is its transpose.
For \(D=J_N/N\), the exact ratio is
\(\prod_{\ell=0}^{r-1}(1-\ell/N)^{-1}\);
its square has first correction \(r(r-1)/N\), as above.
These are symbolic checks on the full quantified domain.
They do not replace the proof of the remainder.
\end{remark}
